\documentclass[a4paper,12pt]{article}
\usepackage[T2A]{fontenc}
\usepackage[utf8]{inputenc}
\usepackage[russian]{babel}
\usepackage{amsmath,amssymb,mathtools}
\usepackage{geometry}
\geometry{top=23mm,bottom=23mm,left=23mm,right=23mm}
\usepackage{microtype}
\usepackage{setspace}
\usepackage{parskip}
\usepackage{array,tabularx,booktabs}
\usepackage{enumitem}
\usepackage{needspace}
\usepackage{xcolor}
\usepackage{tikz}
\usetikzlibrary{arrows.meta,positioning,shapes.geometric,calc}
\usepackage{pgfplots}
\pgfplotsset{compat=1.18}
\usepackage{fancyhdr}
\usepackage{hyperref}
\definecolor{accent}{HTML}{008080}
\definecolor{darkaccent}{HTML}{005757}
\definecolor{textgray}{HTML}{343B3D}
\definecolor{muted}{HTML}{647476}
\hypersetup{colorlinks=true,linkcolor=darkaccent,pdfauthor={Лёвин Артём Александрович},pdftitle={Чек-лист. Построение графиков линейных функций}}
\setstretch{1.20}
\setlength{\parindent}{0pt}
\setlength{\parskip}{0.64em}
\setlist[enumerate]{leftmargin=*,itemsep=.68em,topsep=.5em}
\setlist[itemize]{leftmargin=*,itemsep=.25em,topsep=.4em}
\pagestyle{fancy}
\fancyhf{}
\fancyhead[L]{\small\color{muted}Графики линейных функций}
\fancyhead[R]{\small\color{muted}Кирилл Зиновьев}
\fancyfoot[C]{\small\color{muted}\thepage}
\renewcommand{\headrulewidth}{0.35pt}
\renewcommand{\headrule}{\hbox to\headwidth{\color{accent}\leaders\hrule height \headrulewidth\hfill}}
\setlength{\headheight}{26pt}
\newcommand{\secname}[1]{\needspace{6\baselineskip}\vspace{.85em}\noindent{\Large\color{darkaccent}#1}\par\vspace{-.42em}{\color{accent}\rule{\linewidth}{.45pt}}\vspace{.15em}}
\newcommand{\subname}[1]{\needspace{3\baselineskip}{\color{darkaccent}\bfseries #1}\par\vspace{-.32em}}
\newcommand{\idea}[1]{\noindent{\color{darkaccent}\bfseries Ключевая идея.}\enspace #1}
\newcommand{\warn}[1]{\noindent{\color{darkaccent}\bfseries Проверьте.}\enspace #1}
\newcommand{\smallspace}{\vspace{.35em}}
\begin{document}
\begin{titlepage}
\thispagestyle{empty}
\begin{center}
\vspace*{22mm}
{\large\color{darkaccent} МАТЕМАТИКА}\par
\vspace{24mm}
{\LARGE\color{textgray} Чек-лист}\par
\vspace{15mm}
{\LARGE\color{accent} Построение графиков\\[2mm] линейных функций}\par
\vspace{19mm}
{\large\color{textgray} 7 класс}\par
\vfill
{\large\color{textgray} Кирилл Зиновьев}\par
\vspace{14mm}
{\large\color{textgray} 7 ноября 2026 г.}\par
\vspace{16mm}
{\small\color{textgray} Лёвин Артём Александрович\\[2mm]
эксклюзивно для Кирилла Зиновьева}\par
\vspace{20mm}
\begin{tikzpicture}[x=1cm,y=1cm]
  \draw[accent,line width=.9pt] (0,0) .. controls (1.3,.65) and (2.4,-.62) .. (3.9,-.09)
    .. controls (5.2,.55) and (6.5,-.55) .. (8,0);
\end{tikzpicture}
\end{center}
\end{titlepage}
\setcounter{page}{1}

\secname{1. Что показывает линейная функция}
Линейная функция задаётся формулой
\[
\boxed{y=kx+b,}
\]
где \(x\) --- выбранное значение аргумента, \(y\) --- рассчитанное значение функции, \(k\) и \(b\) --- числа.

График линейной функции --- \textbf{прямая}. Для её построения достаточно \textbf{двух различных точек}.

\begin{itemize}
  \item Если \(k>0\), с ростом \(x\) значение \(y\) возрастает: прямая идёт вверх слева направо.
  \item Если \(k<0\), с ростом \(x\) значение \(y\) убывает: прямая идёт вниз слева направо.
  \item Число \(b\) показывает, где график пересекает ось \(Oy\): \(\boldsymbol{(0;b)}\).
\end{itemize}

\idea{В формуле \(y=2x+3\) число \(2\) умножается на \(x\), а \(3\) прибавляется после умножения. Слагаемые \(2x\) и \(3\) складывать как подобные нельзя.}

\secname{2. Как читать координаты точки}
Точка записывается так: \(A(x;y)\). \textbf{Сначала} указана координата по горизонтальной оси \(Ox\), \textbf{затем} --- по вертикальной оси \(Oy\).

\begin{center}
\begin{tabular}{@{}p{.23\linewidth}p{.67\linewidth}@{}}
\toprule
Точка & Как найти на координатной плоскости\\
\midrule
\(A(0;3)\) & По \(x\) остаёмся на месте; по \(y\) поднимаемся на 3.\\
\(B(2;-1)\) & На 2 вправо, затем на 1 вниз.\\
\(C(-3;5)\) & На 3 влево, затем на 5 вверх.\\
\(D(0;-2)\) & По \(x\) не перемещаемся; по \(y\) опускаемся на 2.\\
\bottomrule
\end{tabular}
\end{center}

\warn{Точки \((0;1)\) и \((1;0)\) разные. При \(x=0\) точка лежит на оси \(Oy\), при \(y=0\) --- на оси \(Ox\).}

\secname{3. Строим прямую по двум точкам}
\subname{Пример: \(y=2x+3\)}
\textbf{Шаг 1.} Самостоятельно выбираем значения \(x\), например \(0\) и \(2\).\\
\textbf{Шаг 2.} Подставляем их в формулу:
\[
\begin{aligned}
x=0&:\quad y=2\cdot0+3=3,\\
x=2&:\quad y=2\cdot2+3=7.
\end{aligned}
\]
\textbf{Шаг 3.} Составляем таблицу:
\[
\begin{array}{c|cc}
 & A&B\\ \hline
x&0&2\\ y&3&7
\end{array}
\]
\textbf{Шаг 4.} Отмечаем \(A(0;3)\) и \(B(2;7)\), проводим через них \textbf{всю прямую}, продолжая её за обе точки.

\begin{center}
\begin{tikzpicture}
\begin{axis}[
 width=11.2cm,height=7.0cm,
 xmin=-2.5,xmax=4.2,ymin=-2.4,ymax=9.3,
 axis lines=middle,
 xlabel={$x$},ylabel={$y$},
 xtick={-2,-1,0,1,2,3,4}, ytick={-2,0,2,3,4,6,7,8},
 grid=both,major grid style={line width=.15pt,draw=gray!35},
 minor tick num=0, tick label style={font=\scriptsize},
 xlabel style={at={(axis description cs:1.02,.18)},anchor=west},
 ylabel style={at={(axis description cs:.42,1.02)},anchor=south},
 clip=false
]
\addplot[very thick,accent,domain=-2.2:3.1,samples=2] {2*x+3};
\addplot[only marks,mark=*,mark size=2.3pt,accent] coordinates {(0,3) (2,7)};
\node[anchor=west,font=\small,fill=white,inner sep=1pt] at (axis cs:0.13,3.25) {$A(0;3)$};
\node[anchor=west,font=\small,fill=white,inner sep=1pt] at (axis cs:2.12,7.05) {$B(2;7)$};
\end{axis}
\end{tikzpicture}
\end{center}

\secname{4. Выбираем удобные значения \(x\)}
\idea{Разрешается подставлять любые удобные \(x\). Особенно полезно выбирать значения, при которых \(y\) получается целым числом.}

\subname{Пример: \(y=-0{,}8x+1\)}
Если взять \(x=0\) и \(x=10\), получим:
\[
\begin{aligned}
y(0)&=-0{,}8\cdot0+1=1,\\
y(10)&=-0{,}8\cdot10+1=-8+1=-7.
\end{aligned}
\]
Точки \(A(0;1)\) и \(B(10;-7)\) имеют целые координаты --- строить их удобно.

При \(x=2\) арифметика требует большего внимания:
\[
y(2)=-0{,}8\cdot2+1=-1{,}6+1=\boxed{-0{,}6}.
\]
\warn{Сумма \(-1{,}6+1\) отрицательна. Минус при вычислении терять нельзя.}

\subname{Пример: \(y=0{,}2x-1\)}
Для графика удобно выбрать \(x=0\) и \(x=5\):
\[
y(0)=-1,\qquad y(5)=0{,}2\cdot5-1=0.
\]
Итак, график проходит через \((0;-1)\) и \((5;0)\). Первая точка лежит на оси \(Oy\), вторая --- на оси \(Ox\).

\secname{5. Масштаб: когда он особенно важен}
Если получаются дробные значения, можно либо подобрать другие \(x\), либо взять более подробный масштаб.

\begin{itemize}
\item Масштаб выбирается \textbf{отдельно для каждой оси}, но в пределах одной оси должен оставаться равномерным.
\item Отметки \(0{,}2\), \(0{,}4\), \(0{,}6\) откладываются через \textbf{равные расстояния}.
\item Все важные деления подпишите. Точку \((2;-0{,}6)\) ставьте \textbf{ниже} оси \(Ox\).
\end{itemize}

\secname{6. Особый случай: \(y=b\)}
Когда в формуле отсутствует видимый \(x\), можно написать \(y=0\cdot x+b\).

\[
\begin{aligned}
y=-2&\quad\Rightarrow\quad (0;-2),\ (3;-2);\\
y=3&\quad\Rightarrow\quad (0;3),\ (5;3).
\end{aligned}
\]
При любом \(x\) значение \(y\) одно и то же. Получается \textbf{горизонтальная прямая}, параллельная оси \(Ox\).

\idea{Чтобы построить \(y=-2\), достаточно найти на оси \(Oy\) значение \(-2\) и провести через него горизонтальную прямую. Если \(y=0\), график совпадает с осью \(Ox\).}

\secname{7. Частые ошибки и самопроверка}
\begin{itemize}
\item \(2x+3\) нельзя упростить до \(5x\): слагаемые имеют разную буквенную часть.
\item Сначала выбирают \(x\), затем \textbf{вычисляют} \(y\) по формуле.
\item Первая координата отвечает за горизонталь, вторая --- за вертикаль.
\item Отрицательная координата \(y\) означает движение \textbf{вниз}.
\item График линейной функции продолжается в обе стороны за выбранные точки.
\item Прямая проходит через начало координат \((0;0)\) только при \(b=0\).
\end{itemize}

\clearpage
\thispagestyle{fancy}
\begin{center}
{\Large\color{darkaccent}Алгоритм построения линейной функции}\par\vspace{8mm}
\end{center}
\begin{center}
\begin{tikzpicture}[
 font=\small,
 >=Stealth,
 action/.style={draw=accent,rounded corners=3pt,align=center,text width=4.6cm,minimum height=1.14cm,inner sep=7pt},
 condition/.style={diamond,aspect=2.0,draw=accent,align=center,text width=2.45cm,inner sep=2pt,minimum height=1.5cm},
 edge/.style={->,line width=.9pt,draw=darkaccent}
]
\node[action] (a) at (0,0) {Прочитайте формулу\\\(y=kx+b\)};
\node[condition] (q) at (0,-2.25) {\(k=0\)?};
\node[action] (yes) at (-3.25,-5.00) {Отметьте \((0;b)\).\\Проведите горизонталь \(y=b\)};
\node[action] (no) at (3.25,-5.00) {Выберите \(x_1=0\) и удобное \(x_2\ne0\)};
\node[action] (calc) at (3.25,-7.2) {Вычислите \(y_1,y_2\).\\Отметьте обе точки};
\node[action] (line) at (3.25,-9.45) {Проведите прямую\\через обе точки};
\node[action,text width=7.0cm] (done) at (0,-12.0) {Проверьте координаты, знаки и масштаб};
\draw[edge] (a) -- (q);
\draw[edge] (q.west) -| node[above,pos=.25]{\color{darkaccent}да} (yes.north);
\draw[edge] (q.east) -| node[above,pos=.25]{\color{darkaccent}нет} (no.north);
\draw[edge] (no) -- (calc);
\draw[edge] (calc) -- (line);
\draw[edge] (line.south) |- (done.east);
\draw[edge] (yes.south) |- (done.west);
\end{tikzpicture}
\end{center}
\clearpage

\secname{Тренировочный блок}
\textit{Путь А \(\cdot\) 10 заданий \(\cdot\) от простого к сложному.}

\smallspace
При построении графика составляйте таблицу минимум для двух точек. Рисунки и расчёты выполните в тетради; ответы сверяйте после самостоятельной работы.

\subname{Базовый уровень}
\begin{enumerate}[label=\textbf{\arabic*.},start=1]
\item Из точек \(A(-3;0)\), \(B(0;-3)\), \(C(0;3)\) выберите ту, у которой \(x=0\), а \(y=-3\). Укажите, на какой оси она находится.

\item Для функции \(y=3x-2\) вычислите \(y\) при \(x=0\) и \(x=2\). Запишите координаты полученных точек.

\item Постройте график \(y=-2x+1\), выбрав \(x=-1\) и \(x=2\). Покажите в таблице обе пары координат.
\end{enumerate}

\subname{Средний уровень}
\begin{enumerate}[label=\textbf{\arabic*.},start=4]
\item Постройте график \(y=0{,}5x+1\), подобрав два значения \(x\), при которых значения \(y\) будут целыми.

\item Постройте график \(y=-0{,}4x+2\), используя \(x=0\) и \(x=5\). Укажите точку его пересечения с осью \(Ox\).

\item Для функции \(y=-0{,}3x+1\) найдите \(y\) при \(x=0\), \(x=10\) и \(x=2\). Для построения графика выберите две точки с целыми координатами.

\item Постройте прямые \(y=-4\) и \(y=2\). Объясните, почему обе параллельны оси \(Ox\).
\end{enumerate}

\subname{Повышенный уровень}
\begin{enumerate}[label=\textbf{\arabic*.},start=8]
\item Какие из графиков проходят через начало координат: \(y=-3x\), \(y=2x+5\), \(y=0\), \(y=-4\)? Объясните выбор через значение \(b\).

\item Для функции \(y=-0{,}6x+4\) проверьте точки \(A(0;4)\), \(B(5;1)\), \(C(10;-1)\). Какие из них принадлежат графику? Постройте прямую по двум подходящим точкам.

\item На одной координатной плоскости постройте графики \(y=0{,}25x-2\) и \(y=-3\). Подберите удобные точки, отметьте пересечения с осью \(Oy\) и укажите горизонтальную прямую.
\end{enumerate}

\clearpage
\secname{Ответы к тренировочному блоку}
\renewcommand{\arraystretch}{1.40}
\begin{tabularx}{\linewidth}{@{}p{1.2cm}X@{}}
\toprule
\textbf{№} & \textbf{Ответ / ориентир для проверки}\\
\midrule
1 & \(B(0;-3)\), ось \(Oy\).\\
2 & \((0;-2)\), \((2;4)\).\\
3 & \((-1;3)\), \((2;-3)\).\\
4 & Например, \((0;1)\), \((4;3)\).\\
5 & \((0;2)\), \((5;0)\); пересечение с \(Ox\): \((5;0)\).\\
6 & \(y(0)=1\), \(y(10)=-2\), \(y(2)=0{,}4\); для графика: \((0;1)\), \((10;-2)\).\\
7 & Горизонтальные прямые через \((0;-4)\) и \((0;2)\); \(k=0\).\\
8 & \(y=-3x\) и \(y=0\), поскольку \(b=0\).\\
9 & \(A\) и \(B\) принадлежат; \(C\) --- нет (при \(x=10\) имеем \(y=-2\)).\\
10 & \(y=0{,}25x-2\): \((0;-2)\), \((8;0)\).\\
& \(y=-3\): \((0;-3)\), \((4;-3)\); горизонтальна вторая.\\
\bottomrule
\end{tabularx}

\vspace{11mm}
\subname{Итоговая памятка}
\begin{center}
\begin{tabularx}{\linewidth}{@{}p{.31\linewidth}X@{}}
\toprule
\textbf{Сигнал в задаче} & \textbf{Что делать}\\
\midrule
Дана формула \(y=kx+b\) & Выберите два удобных \(x\), рассчитайте \(y\), отметьте две точки.\\
Десятичный коэффициент & Попробуйте \(x=5\) или \(x=10\), чтобы избежать дробных координат.\\
В формуле \(y=b\) нет \(x\) & Постройте горизонтальную прямую через \((0;b)\).\\
Точка \((0;b)\) & Это пересечение прямой с осью \(Oy\).\\
\bottomrule
\end{tabularx}
\end{center}
\vfill
\begin{center}
{\small\color{muted}Лёвин Артём Александрович \(\cdot\) эксклюзивно для Кирилла Зиновьева}
\end{center}
\end{document}
